\chapter{Historical Case Studies}\label{ch:history}

\begin{objectives}
\begin{itemize}
\item connect each design concern in this tutorial to a real event that made it famous;
\item recompute the key numbers of three numerical failures and check them against the official accounts;
\item extract from each case one concrete rule for \texttt{score}.
\end{itemize}
{\small\color{muted}\textbf{Chapter references:} Kegel~\cite{kegel1999c10k}; Banga et al.~\cite{banga1999scalable}; Lemon~\cite{lemon2001kqueue}; Axboe~\cite{axboe2019iouring}; Lowenstein~\cite{lowenstein2000genius}; BCBS~\cite{bcbs2009revisions,bcbs2019frtb}; GAO~\cite{gao1992patriot}; Skeel~\cite{skeel1992roundoff}; Lions et al.~\cite{lions1996ariane}.}
\end{objectives}

Each case below is told with the numbers we could recompute or source to a primary account; where accounts differ or a
figure could not be checked, we say so and state less.

\section{C10K, prefork and the event loop}\label{sec:hist-c10k}
In 1999 Dan Kegel posted ``The C10K problem''~\cite{kegel1999c10k}: hardware could plausibly serve ten thousand
concurrent clients, but the dominant software designs could not. The widespread Apache configuration of the time, the
\emph{prefork} model, ran one process per connection, the extreme version of \cref{def:tpcmodel}. Kegel's page surveyed the
alternatives: non-blocking I/O with \cc{select}/\cc{poll}, signal-driven I/O, kernel event queues, and asynchronous I/O.
Banga, Mogul and Druschel showed experimentally that the cost of \cc{select} grew with the number of descriptors and
proposed a stateful, explicit event interface~\cite{banga1999scalable}, the design idea that \cref{thm:epoll} formalises.
FreeBSD shipped \cc{kqueue} in 2000 and Jonathan Lemon described it at USENIX in 2001~\cite{lemon2001kqueue}; Linux gained
\cc{epoll} in the 2.5 development series and in stable 2.6 kernels. nginx, first publicly released in 2004, was built on
such event loops from the start. Fifteen years later, io\_uring (Linux 5.1, 2019)~\cite{axboe2019iouring} moved the frontier
from ``one system call per wait'' to ``shared rings, few system calls at all''.

\textbf{Rule for \texttt{score}.} Measure before choosing, and choose behind an interface: \cref{fig:pollbench} reproduces
the 1999 observation on a 2020s laptop in a hundred lines of C++, and \cc{EventLoop} keeps the choice reversible.

\section{Long-Term Capital Management, 1998}\label{sec:hist-ltcm}
LTCM was founded in 1994 by John Meriwether; its partners included Myron Scholes and Robert Merton, who shared the 1997
Nobel memorial prize in economics. Its strategies were highly leveraged convergence trades whose risk models were
estimated on historical data from relatively calm markets. After Russia defaulted on its domestic debt in August 1998,
investors fled to the most liquid assets, spreads that the models expected to narrow widened together, and
positions that looked diversified lost money at the same time. In September 1998 the Federal Reserve Bank of New York
organised a recapitalisation of about \$3.6 billion by a consortium of private financial institutions~\cite{lowenstein2000genius}.

\textbf{Rule for \texttt{score}.} Historical simulation can only see the past it is given. A 503-day window that happens to
be calm (\cref{sec:mockvar}) says nothing about a regime it has not seen, which is why \cref{ex:kupiec} asks for out-of-sample
backtests and why stressed windows matter.

\section{The 2008 crisis and the move to Expected Shortfall}\label{sec:hist-2008}
Trading-book capital computed from 99\% VaR proved too low in 2007--2008, for two reasons this tutorial has made precise:
VaR ignores the size of losses beyond its quantile (\cref{sec:es}), and models fitted to the preceding calm years
under-predicted the tail. The Basel Committee responded first with ``Basel 2.5'' in 2009, which added a \emph{stressed}
VaR calibrated to a period of significant stress~\cite{bcbs2009revisions}, and then with the Fundamental Review of the
Trading Book, which replaced VaR with 97.5\% Expected Shortfall~\cite{bcbs2019frtb}.

\textbf{Rule for \texttt{score}.} Report CVaR beside VaR, always. The repository already computes both from the same sorted tail.

\section{The Patriot battery at Dhahran, 1991}\label{sec:hist-patriot}
On 25 February 1991 a Patriot air-defence battery at Dhahran, Saudi Arabia, failed to engage an incoming Scud missile,
which struck a US Army barracks and killed 28 American soldiers. The US General Accounting Office traced the failure to
the system's timekeeping~\cite{gao1992patriot}. The system counted time as an integer number of tenths of a second and
converted it to seconds by multiplying by $1/10$ held in a 24-bit fixed-point register. Because $1/10$ has no finite binary
expansion (\cref{lem:dyadic}), the stored constant was truncated, and the error grew linearly with uptime. The battery had
been running for about 100 hours.

\begin{example}[Recomputing the drift]\label[example]{ex:patriot}
Skeel's analysis~\cite{skeel1992roundoff} takes the stored value to be $1/10$ chopped after 23 binary places:
\[ 0.00011001100110011001100_2=\tfrac{209715}{2097152},\qquad \tfrac1{10}-\tfrac{209715}{2097152}=9.5367431640625\times10^{-8}. \]
After $h$ hours the clock has ticked $36{,}000h$ times, so the time error is $36{,}000h\times9.5367\times10^{-8}$\,s.
\begin{center}
\small
\begin{tabular}{@{}rrrrrrr@{}}
\toprule
\textcolor{muted}{hours} & 1 & 8 & 20 & 48 & 72 & 100\\
\midrule
\textcolor{muted}{drift (s)} & 0.0034 & 0.0275 & 0.0687 & 0.1648 & 0.2472 & 0.3433\\
\bottomrule
\end{tabular}
\end{center}
These values (\file{verify/hist.py}) match the GAO report's table, including $0.3433$\,s at 100 hours, which GAO
translates into a shift of 687 metres in the range gate where the radar looked for the target. Chopping after 24 binary places
instead would give an error of $3.58\times10^{-8}$ per tick and $0.1287$\,s at 100 hours, which does \emph{not} match;
the published numbers pin down the 23-place interpretation.
\end{example}

\begin{figure}[htbp]
\centering
\begin{tikzpicture}
\begin{axis}[darkaxis,width=0.82\linewidth,height=5.4cm,xlabel={hours of continuous operation},ylabel={clock error (s)},
  grid=major,xmin=0,xmax=105,ymin=0,ymax=0.38]
\addplot[red,thick,mark=*,mark size=1.8pt] table[x=h,y=drift]{data/patriot.dat};
\node[text=red,font=\footnotesize,anchor=south east] at (axis cs:100,0.3433) {0.3433 s at 100 h};
\end{axis}
\end{tikzpicture}
\caption{Patriot clock drift from a truncated binary $1/10$, recomputed exactly with rational arithmetic (pgfplots).}
\label{fig:patriot}
\end{figure}

According to GAO, Israeli users had reported the loss of accuracy after long runs, and a software fix was issued but reached
Dhahran on 26 February, the day after the attack~\cite{gao1992patriot}.

\textbf{Rule for \texttt{score}.} An error of one part in ten million per step is harmless once and fatal when multiplied by
millions of steps. The same arithmetic applies to an epsilon guard: \cref{thm:epsfloor} bounds the error as a function of $n$.

\section{The Vancouver Stock Exchange index, 1982--83}\label{sec:hist-vse}
The Vancouver Stock Exchange launched a new index in January 1982 at a base value of 1000. It was recomputed after every
trade, reportedly a few thousand times a day, and each new value was \emph{truncated}, not rounded, to three decimal
places. By late 1983 the index stood at roughly half the value it should have had; the figures usually quoted are about
524.8 before correction and about 1098.9 after recomputation.

The mechanism is easy to check. Truncation to three decimals discards on average half of the last unit, $0.0005$ points,
on every update, always in the same direction, while rounding errors cancel on average. \Cref{fig:vse} simulates
$1{,}386{,}000$ updates ($3000$ per day for 22 months of 21 trading days) of a driftless random index with tiny moves.

\begin{figure}[htbp]
\centering
\begin{tikzpicture}
\begin{axis}[darkaxis,width=0.9\linewidth,height=6.0cm,xlabel={number of recomputations},ylabel={index value},
  grid=major,legend pos=south west,ymin=0,ymax=1150,scaled x ticks=base 10:-6,xtick scale label code/.code={$\times10^{6}$},
  every x tick scale label/.style={at={(1,0)},anchor=north west,font=\footnotesize,text=muted}]
\addplot[grn,thick] table[x=u,y=exact]{data/vse.dat}; \addlegendentry{exact (no rounding)}
\addplot[acc,thick,dashed] table[x=u,y=round]{data/vse.dat}; \addlegendentry{rounded to 3 decimals}
\addplot[red,thick] table[x=u,y=trunc]{data/vse.dat}; \addlegendentry{truncated to 3 decimals}
\end{axis}
\end{tikzpicture}
\caption{Simulation of truncation bias (\file{verify/make_plot_data.py}, seed 1982; illustrative, not historical data).
After $1.386\times10^6$ updates the exact path ends at 930.18, the rounded one at 930.37, and the truncated one at 275.81:
a mean loss of $0.00047$ points per update, close to the $0.0005$ predicted.}
\label{fig:vse}
\end{figure}

\textbf{Rule for \texttt{score}.} A biased rounding mode is worse than a large unbiased one, because bias accumulates
linearly while unbiased error grows like $\sqrt{n}$. The cast \cc{static_cast<std::size_t>} in \cref{lst:epsfix} is a
truncation, and it was a truncation that turned $0.99999999999999978$ into $0$.

\section{Ariane 5 Flight 501, 1996}\label{sec:hist-ariane}
On 4 June 1996, about 37 seconds after lift-off, the first Ariane~5 veered off course and broke up. The Inquiry Board
chaired by Jacques-Louis Lions reported on 19 July 1996~\cite{lions1996ariane}. The inertial reference software, reused from
Ariane~4, converted a 64-bit floating-point value related to horizontal velocity (the ``horizontal bias'') to a 16-bit
signed integer. Ariane~5's faster trajectory produced a value larger than $2^{15}-1=32{,}767$, the conversion raised an
operand error that was not handled, and both the active and the backup inertial units, running the same code, shut down.
The flight computer then interpreted diagnostic data as flight data and commanded full nozzle deflections. The function
that failed served no purpose after lift-off; it kept running for compatibility with Ariane~4's requirements.

\textbf{Rule for \texttt{score}.} Every narrowing conversion is a claim about the range of its input. The \cc{static_cast}
in \cc{historical_tail} narrows a \cc{double} to \cc{std::size_t} and is safe only because the function first checks
$0<c<1$; the length prefix in \cref{ch:framing} narrows a \cc{size_t} to \cc{uint32_t}, and \cref{ex:maxframe} asks what
guards it.

\section{Synthesis}
\begin{center}
\footnotesize
\begin{tabularx}{\linewidth}{@{}>{\leavevmode\bfseries\color{yel}}p{2.4cm} >{\raggedright\arraybackslash}X >{\raggedright\arraybackslash}X >{\raggedright\arraybackslash}p{2.1cm}@{}}
\toprule
\textcolor{muted}{Case} & \textcolor{muted}{What failed} & \textcolor{muted}{Lesson applied in \texttt{score}} & \textcolor{muted}{See}\\
\midrule
C10K (1999) & per-connection processes and $O(n)$ readiness checks & concurrency model behind \cc{EventLoop}; measure & \cref{ch:conc}\\
LTCM (1998) & models fitted to calm history & out-of-sample backtests, stressed windows & \cref{ch:risk}\\
2008 / FRTB & VaR blind beyond its quantile & report CVaR with VaR & \cref{sec:es}\\
Patriot (1991) & truncated binary $1/10$, error $\propto$ uptime & error bounds as a function of $n$ & \cref{thm:epsfloor}\\
VSE (1982--83) & truncation bias accumulating & avoid biased casts on boundary values & \cref{lst:epsfix}\\
Ariane 5 (1996) & unguarded narrowing conversion & range checks before every narrowing & \cref{ch:framing,ch:float}\\
\bottomrule
\end{tabularx}
\end{center}

\begin{exercises}
\ex{ex:patriotfix}{\lvB} Give two fixes for the Patriot clock that need no extra precision: one that changes the
representation of time, one that changes the computation. Which of them corresponds to the basis-point fix of
\cref{lst:bp}?

\ex{ex:vsebias}{\lvB} Prove that if the true fractional part (in units of $0.001$) of each update is uniform on $[0,1)$,
truncation loses $0.0005$ points per update in expectation, while round-half-to-even loses $0$. How large must the
per-update moves be for the uniform assumption to be reasonable?

\ex{ex:ariane}{\lvC} The Ariane~4 code was not wrong for Ariane~4. State the precondition the conversion relied on as a
formula, and design a test strategy for \texttt{score} that would detect when a reused component's preconditions no longer
hold (think of \cc{historical_tail} reused by a future intraday store with $n$ in the millions).
\end{exercises}
